\subsection{The drafts}
\label{sec:11.4.1}
Two bills in the current Congress answer the first question. The Abolish ICE Act would cut off federal funds for the functions assigned to ICE's director under the 2002 Act, rescind the unobligated balances of what had been made available to its director, and abolish the agency ninety days after enactment \autocite{hr7123}. The Melt ICE Act works through the agency's powers rather than its name \autocite{hr7190}. It releases everyone held in DHS detention within six months, repeals the detention authorities and the 287(g) delegation to state and local police, ends detention and monitoring contracts and then bars all federal funds for detention and monitoring, and bars ICE's operations account from paying for civil arrests, detention, removal or charging documents. Appendix~\ref{sec:C} lists its provisions.

Two further drafts reach one entry of the schedule. The ICE Out of Our Faces Act, S. 3779, would bar ICE and CBP from acquiring or using facial recognition and other biometric identification systems, require deletion of the data collected for them, and let individuals and state attorneys general sue for violations \autocite{Jayapal2026FacesAct}. The Facial Recognition and Biometric Technology Moratorium Act, reintroduced in September 2026, would bar every federal entity from using facial recognition and other biometric technologies until Congress lifts the bar. It would also make federal grants to state and local agencies conditional on their adopting moratoria of their own \autocite{Jayapal2026Moratorium}.

This Congress passed the Secure America Act 214–212 and 52–47 on party lines \autocite{GT2026SecureAmerica}, and none of these bills has a path to enactment in it. They exist as drafts, written before anyone needs them, so that enacting one is a matter of votes and not of drafting.
